\section{Platform Design and Calibration Motivation}

\subsection{Processing Architecture}

The platform was selected according to five criteria: low relative cost, reconfigurable analog routing, hardware-assisted signal processing, reuse of acquisition resources across operating modes, and sufficient conversion resolution for both measurement and internal diagnostics. In acquisition mode, the PSoC conditions and digitizes the geophone output. In calibration mode, the input chain remains unchanged, but the observation path, digital-filter coefficients, and DAC-control firmware are reconfigured to close a temporary mixed-signal loop.

Self-calibration using on-chip resources is not itself new. The PSoC-Stat injects five known current levels with an IDAC and uses the resulting ADC values to estimate the gain and offset of a TIA/ADC signal chain \cite{Lopin2018}. Infineon's AN68403 applies an internal VDAC to a PGA/ADC chain, measures the system offset with grounded inputs, and stores gain and offset corrections in EEPROM \cite{InfineonAN68403}. These methods primarily correct the numerical transfer characteristic observed at the ADC output. The objective here is different: the proposed loop changes the physical reference of every amplification stage so that each intermediate analog node returns to its intended operating point.

The available ADDER made it possible to implement a fixed DC injection as an early workaround. This solution was viable and functional for this particular circuit, but it is not generally transferable to front ends without an explicit summing node. More importantly, it is only a patch: earlier stages continue to drift before the correction point, while subsequent stages introduce new offsets after it.

For the $i$th stage, its slowly varying output can be represented as
\begin{equation}
 y_i^{\mathrm{DC}} = A_i y_{i-1}^{\mathrm{DC}} + b_i + K_i u_i,
 \label{eq:stage}
\end{equation}
where $A_i$ is the stage gain, $b_i$ contains the amplifier and reference errors, $K_i$ is the signed transfer from the trim DAC to the observed node, and $u_i$ is the DAC code. Adjusting only the ADDER reference can force one selected output toward its desired value, but it does not constrain the remaining $y_i^{\mathrm{DC}}$ values. If an earlier stage reaches a rail, the summing correction cannot reconstruct the lost waveform; offsets appearing after the ADDER are also left uncompensated.

The proposed architecture instead chooses $u_i$ independently so that every observed $y_i^{\mathrm{DC}}$ approaches its stage reference. In the differential GEO implementation, all four calibration targets correspond to zero ADC counts.

